\documentclass{article}

\usepackage{amsmath} % for \text
\usepackage{amssymb} % for \mathbb
\usepackage{amsthm} % for proof environment
\usepackage{hyperref} % for \autoref
\usepackage{cleveref} % for \cref
\usepackage{tikz} % for tikzpicture
\usepackage{graphicx} % for \includegraphics
\usepackage{subcaption} % for subfigure
\usepackage{float} % for H in figure
\usepackage{geometry} % for \newgeometry
\usepackage{fullpage} % for fullpage


\geometry{left = 1.8cm, right = 1.8cm}

\begin{document}

\begin{center}
    \textbf{\LARGE{Ecole Normale Supérieure de Paris-Saclay}}\\
\end{center}

\vspace{1cm}

\begin{center}
    \textbf{Rapport TER}
\end{center}

\noindent\rule{\textwidth}{0.6mm}
\bigskip
\begin{center}
    \textbf{\LARGE{TER - Voiture autonomes avec apprentissage par renforcement et lidar}}
\end{center}
\bigskip
\noindent\rule{\textwidth}{0.6mm}

\bigskip
\begin{center}
    \textsc{Miquel Hugo}
\end{center}
\begin{center}
    \textsc{Plus Basile}
\end{center}

\begin{figure}[b]
    \centering
    \begin{minipage}[h]{0.45\textwidth}
        \raggedright
        \includegraphics[width=0.6\textwidth]{Images/Logo-ENS-Paris-Saclay.png}
    \end{minipage}
    \begin{minipage}[h]{0.45\textwidth}
        \raggedleft
        \includegraphics[width=0.6\textwidth]{Images/Logo-Universite-Paris-Saclay.jpg}
    \end{minipage}
\end{figure}


\newpage

\tableofcontents

\newpage

\section{Introduction}

\subsection{Contexte}

Les voitures autonomes sont un sujet de recherche très actif depuis quelques années. 
En effet, elles pourraient révolutionner le monde des transports en permettant de réduire les accidents de la route, 
de diminuer la consommation d'énergie et de réduire les embouteillages. 
Cependant, il reste encore de nombreux défis à relever pour que les voitures autonomes soient utilisées à grande échelle.
En particulier, il est nécessaire de développer des algorithmes d'apprentissage par renforcement qui permettent à une 
voiture autonome d'apprendre à conduire de manière autonome.


\end{document}